% profiles/university-of-connecticut.tex — University of Connecticut
% ============================================================================
% source: https://registrar.uconn.edu/graduation/doctoral-degrees/
%         ("Dissertation Information" / "Dissertation Specifications", the
%         numbered filing requirements #1-#14)
% source: https://registrar.media.uconn.edu/wp-content/uploads/sites/1604/2026/07/Dissertation-Samples-1-4-2-1_TimesNewRoman.docx
%         (the sample the page links as a starting document — "Use these
%         properly formatted, accessible dissertation samples ... as starting
%         documents for guidance"; an Arial twin sits beside it. Uploaded
%         July 2026, the most recent artifact of the set.)
% accessed: 2026-09-13
% checked:  texlib-thesis-profile-round
%
% WHY THE REGISTRAR AND NOT A GRADUATE SCHOOL. UConn's dissertation filing
% requirements are published by the Office of the Registrar, which administers
% degree completion there; grad.uconn.edu carries no competing specifications
% page. This is the institution's own filing requirement document, not a
% department page or a library guide, and it is the one the submission portal
% enforces.
%
% VERIFIED and set here: geometry, spacing, title page, the advisory committee
% page, and a full set of accessibility requirements.
% NOT set: chapter-opening margin (UConn states none) and the typeface (see
% below — UConn "recommends", and adds a font-format rule this class cannot
% satisfy literally).

\thesisinstitution{University of Connecticut}

% "Minimum Margins: The minimum acceptable margins for all pages of the
% dissertation and the abstract are 1 inch on left and 1 inch on the top,
% bottom, and right."
%
% Note "minimum" — wider is in spec, and UConn states no binding allowance.
% Note also that the rule explicitly covers the ABSTRACT, which at UConn is a
% separate unpaginated document ahead of the title page.
\thesissetgeometry{left=1in, right=1in, top=1in, bottom=1in}

% "Spacing: The text of the dissertation should be double-spaced. Long
% quotations, footnotes, appendices, and references may be single-spaced and
% indented."
\thesissetspacing{double}

% --- Typeface: NOT set, and there are two separate reasons ---------------------
% "Font and Point Size: Recommended fonts include Arial, Times New Roman, and
% Helvetica with a point size of either 11 or 12. For accessibility purposes, a
% sans serif font (e.g., Arial, Helvetica) is recommended, unless a specific
% serif font is an industry standard. The font selected must be an embedded,
% True Type (not scalable) font."
%
% The first half is "recommended", twice, so there is nothing to require. The
% second half is a "must" this class cannot satisfy to the letter: lualatex
% embeds OpenType faces as Type0/CID rather than as TrueType. Everything IS
% embedded, which is the substance of the rule, but a filer bounced on the
% font format should know the sentence exists. No font choice available to this
% file would change how lualatex embeds, which is why it is recorded here
% rather than encoded.

% --- Accessibility -----------------------------------------------------------
% UConn is the second school in this round with enforceable accessibility rules
% and the only one that names WCAG. The requirement is dated, and the date is
% in the near future rather than the past:
%
%   "EFFECTIVE for SPRING 2027 GRADUATES:
%    Accessibility: In accordance with the final ruling by the Department of
%    Justice which updated Title II of the Americans with Disabilities Act,
%    your dissertation must meet the Web Content Accessibility Guidelines
%    (WCAG) 2.1."
%
% Spring 2027 is the next graduating class for anyone filing now, so these are
% recorded as requirements rather than as advice. (Note the Title II date here
% is UConn's own; other profiles in this set carry different compliance dates
% for the same rule, legitimately — large public entities and smaller ones have
% different deadlines. Do not reconcile them against each other.)
%
% UConn also restates it in the body of the specifications: "The copy of your
% dissertation that is uploaded must be digitally accessible, as outlined above
% in #1."
\thesisrequiretagging

% THE ONE INFERENCE IN THIS FILE, made deliberately and flagged. UConn's
% enumerated list does not name the document language; what it names is
% conformance to WCAG 2.1 as a whole, and "Language of Page" (SC 3.1.1) is a
% Level A criterion of WCAG 2.1, so it is required by reference. The list is
% introduced with "This includes", which is non-exhaustive on its face.
%
% It is set rather than omitted because the asymmetry favours it: a false
% positive here costs a build-time warning, while a false negative costs a
% filed document that fails the standard UConn names. A reviewer who wants only
% what UConn enumerates should delete this one line; nothing else depends on it.
\thesisrequiredocumentlanguage

% NOT set: \thesisrequirepdfstandard. UConn names no PDF standard anywhere —
% it names WCAG, which is a web content standard, and lists document
% properties. Inventing PDF/A or PDF/UA here would be a claim about what the
% submission portal checks.

\thesisaccessibilityrequirement{WCAG 2.1}
	{In accordance with the final ruling by the Department of Justice which
	 updated Title II of the Americans with Disabilities Act, your dissertation
	 must meet the Web Content Accessibility Guidelines (WCAG) 2.1. Effective
	 for Spring 2027 graduates.}
\thesisaccessibilityrequirement{Document title}
	{Document Title (Word: File > Info or PDF: File > Document Properties).}
\thesisaccessibilityrequirement{Headings}
	{Headings (H1 title of document, H2 main sections, H3 subsections), using
	 Styles in Microsoft Word or a similar feature in other applications. The
	 sample document sets the dissertation title itself in Heading 1.}
\thesisaccessibilityrequirement{Table headings}
	{Ensure tables have a clear, simple structure with properly designated
	 header rows and/or columns.}
\thesisaccessibilityrequirement{Alt text}
	{Alt text or captions for images. Photographs, graphics, and scanned images
	 should be high quality and include a text alternative (Alt Text) that
	 serves the equivalent purpose for each image.}
\thesisaccessibilityrequirement{Colour contrast}
	{Adequate color contrast for text. Use an accessibility checker, such as
	 Microsoft Word or Acrobat Pro, and a color contrast checker, such as
	 WebAIM.}

% --- Profile-local fields -----------------------------------------------------
% UConn's title page lists the author's PRIOR degrees, one per line, under the
% name — "B.A., Smith State College, [year awarded] / M.A., University of
% Connecticut, [year awarded]" — and repeats the abbreviations after the name
% on the committee page ("Jane Mary Doe, B.A., M.A."). Neither has a slot in
% the class's metadata model.
%
% Both default to empty rather than to a placeholder: a student with no prior
% graduate degree legitimately has little or nothing to put here, so an empty
% value is a real answer and the lines simply do not print.
\newcommand{\thesis@uc@priordegrees}{}
\newcommand{\uconnpriordegrees}[1]{\renewcommand{\thesis@uc@priordegrees}{#1}}
\newcommand{\thesis@uc@postnominals}{}
\newcommand{\uconnpostnominals}[1]{\renewcommand{\thesis@uc@postnominals}{#1}}
% "[year of graduation]" on both pages. Its own field because UConn asks for a
% bare year where \graddate usually holds a term and a year.
\newcommand{\thesis@uc@year}{\the\year}
\newcommand{\uconnyear}[1]{\renewcommand{\thesis@uc@year}{#1}}

% --- Title page --------------------------------------------------------------
% The sample's layout. Everything on the page is centred: the title paragraph
% carries no explicit justification of its own but is styled Heading 1, and
% Heading 1 in that document's styles.xml sets w:jc="center". (Read the style,
% not the paragraph — the paragraph alone looks left-aligned.)
%
%   The Title of the Dissertation in Upper and Lower Case Letters
%   Jane Mary Doe
%   B.A., Smith State College, [year awarded]
%   M.A., University of Connecticut, [year awarded]
%   A Dissertation
%   Submitted in Partial Fulfillment of the
%   Requirements for the Degree of
%   Doctor of Philosophy
%   at the
%   University of Connecticut
%   [year of graduation]
%   i (may be suppressed)
%
% "Upper and Lower Case Letters" is title case, so \@title is set as typed and
% NOT uppercased. The page number may be suppressed, and is here: "Pagination:
% ... The title page number may be suppressed."
%
% ONE THING THE SOURCES DO NOT PRESCRIBE: where the block sits vertically. The
% specifications say nothing about it, and the sample is a Word document whose
% spacing comes from empty paragraphs rather than from a stated measurement, so
% reading an exact figure off it would be reading more than it says. The block
% is centred vertically here, which is the conservative choice; if a reviewer
% has a UConn filing to compare against and it sits higher, this is the line to
% change.
\renewcommand{\thesistitlepage}{%
	\loadgeometry{nopagenumber}%
	\begin{titlepage}%
		\thispagestyle{empty}\singlespacing
		\vspace*{\fill}
		{\centering
			\thesis@tagtitle{Title}{\@title}\par
			\vspace{3\baselineskip}
			\@author\par
			\ifx\thesis@uc@priordegrees\@empty\else
				\vspace{2\baselineskip}
				\thesis@uc@priordegrees\par
			\fi
			\vspace{3\baselineskip}
			A \thesis@DocNoun\par
			Submitted in Partial Fulfillment of the\par
			Requirements for the Degree of\par
			\thesis@degree\par
			at the\par
			University of Connecticut\par
			\thesis@uc@year\par
		\par}
		\vspace*{\fill}
	\end{titlepage}%
	\loadgeometry{normal}%
}

% --- Advisory committee page --------------------------------------------------
% UConn calls it the advisory committee page and gives it a place in the
% sequence: "The appropriate order of the major sections of the dissertation
% follows: the abstract, the title page, the copyright page (optional), the
% advisory committee page, acknowledgments, table of contents, the text,
% appendices (if any), and the references."
%
% The sample's layout, headed by the degree rather than by the word "Approval":
%
%   Doctor of Philosophy Dissertation
%   The Title of the Dissertation in Upper and Lower Case Letters
%   Presented by
%   Jane Mary Doe, B.A., M.A.
%   Approved by
%   Major Advisor: John A. Jones
%   Associate Advisor:  Mary J. Smith
%   Associate Advisor:  James S. Williams
%   University of Connecticut
%   [year of graduation]
%
% NO SIGNATURE RULES, and that is what the sample shows: typed names on the
% same line as their role. UConn does collect real signatures, but on the
% approval FORMS rather than on this leaf — "Original, scanned, or electronic
% signatures are accepted", and "All members of a student's advisory committee
% must provide an original signature ... however, signatures may be on
% different pages or come from multiple faculty emails". So the page records
% who approved; the forms carry the signing.
%
% This page IS numbered (it sits inside the Roman-numbered preliminary pages),
% so unlike the class's neutral approval page it does not set
% \thispagestyle{empty}. The document still owns \pagenumbering{roman}.
\renewcommand{\thesisapprovalpage}{%
	\clearpage
	\loadgeometry{normal}\singlespacing
	{\centering
		\thesis@degree{} \thesis@DocNoun\par
		\vspace{2\baselineskip}
		\@title\par
		\vspace{3\baselineskip}
		Presented by\par
		\vspace{\baselineskip}
		\@author\ifx\thesis@uc@postnominals\@empty\else
			, \thesis@uc@postnominals
		\fi\par
		\vspace{3\baselineskip}
		Approved by\par
		\vspace{\baselineskip}
		% Role and name on one line, no rule. The class default
		% \thesis@memberline draws a signature rule and puts the role AFTER the
		% name; UConn's sample does neither.
		\begingroup
		\renewcommand{\thesis@memberline}[2]{%
			\ifx\relax##2\relax##1\else##2: ##1\fi\par
		}%
		Major Advisor: \thesis@advisor\par
		\thesis@committee
		\endgroup
		\vspace{3\baselineskip}
		University of Connecticut\par
		\thesis@uc@year\par
	\par}
	\clearpage
}
